% section17.tex
% §3: The Three-Ray Eigenvalue Theorem
% For inclusion in the NT paper (amsart class)
% All results verified computationally in SageMath.
% Author: Marvin L. Gentry

\section{The Three-Ray Eigenvalue Theorem}
\label{sec:three-ray}

\subsection{Setup}
\label{subsec:three-ray-setup}

By Theorem~\ref{thm:frobenius}, the $\chi_5^{(31)}$-twist of $f_{283}$
produces a Bianchi newform $g$ of level $\mathfrak{n}_{8773}$ over
$\mathbb{Q}(\sqrt{-3})$, with Hecke eigenvalue field $\mathbb{Q}(\sqrt{5})$.
We now determine the eigenvalues of $g$ explicitly.

Fix a primitive root $\mathfrak{g} = 3$ modulo 31
(since $3^1 = 3$, $3^2 = 9$, $3^3 = 27$, $3^5 = 243 \equiv 26$,
$3^{10} \equiv 25 \equiv -6$, $3^{15} \equiv -1$, $3^{30} \equiv 1$).
For each prime $p \neq 2, 3, 31, 283$, define
\[
  k(p) \;:=\; \mathrm{dlog}_{\mathfrak{g}}(p \bmod 31) \bmod 5
  \;\in\; \{0,1,2,3,4\},
\]
where $\mathrm{dlog}_{\mathfrak{g}}$ denotes the discrete logarithm base
$\mathfrak{g} = 3$ in $(\mathbb{Z}/31\mathbb{Z})^\times$.

The character $\chi_5 = \chi_5^{(31)}$ of conductor 31 and order 5 is
the unique primitive character with
\[
  \chi_5(3) = \zeta_5,
\]
where $\zeta_5 = e^{2\pi i/5}$. Then $\chi_5(p) = \zeta_5^{k(p)}$.

The fifth-roots of unity satisfy the identities
\[
  \zeta_5 + \zeta_5^{-1} = \frac{\sqrt{5}-1}{2} = \frac{1}{\varphi},
  \qquad
  \zeta_5^2 + \zeta_5^{-2} = \frac{-\sqrt{5}-1}{2} = -\varphi,
\]
where $\varphi = (1+\sqrt{5})/2$ is the golden ratio.
Define the \emph{character sums}
\begin{equation}
\label{eq:Sk}
  S_k \;:=\; \zeta_5^k + \zeta_5^{-k} \;\in\; \mathbb{Q}(\sqrt{5}),
  \qquad k \in \{0,1,2,3,4\},
\end{equation}
which take exactly three distinct values:
\[
  S_0 = 2, \qquad
  S_1 = S_4 = \tfrac{1}{\varphi} = \tfrac{\sqrt{5}-1}{2}, \qquad
  S_2 = S_3 = -\varphi = \tfrac{-\sqrt{5}-1}{2}.
\]

\subsection{The main theorem}
\label{subsec:three-ray-thm}

\begin{theorem}[Three-Ray Eigenvalue Theorem]
\label{thm:three-ray}
Let $g$ be the $\mathbb{Q}(\sqrt{5})$-Bianchi newform at level $8773 = 283
\cdot 31$ produced by Theorem~\ref{thm:frobenius}. For every prime
$p \neq 2, 3, 31, 283$, the Hecke eigenvalue of $g$ at $p$ is
\begin{equation}
\label{eq:three-ray}
  a_p(g) \;=\; a_p(f_{11}) \cdot S_{k(p)},
\end{equation}
where $k(p) = \mathrm{dlog}_3(p \bmod 31) \bmod 5$ and $S_k$ is as in
\eqref{eq:Sk}.

Equivalently, writing $a_p(g) = A + B\sqrt{5}$ with
$A, B \in \tfrac{1}{2}\mathbb{Z}$:
\begin{equation}
\label{eq:three-ray-coords}
  (A, B) \;=\; \begin{cases}
    (2\,a_p,\; 0)          & \text{if } k(p) = 0, \\[4pt]
    (-\tfrac{a_p}{2},\; +\tfrac{a_p}{2}) & \text{if } k(p) \in \{1,4\}, \\[4pt]
    (-\tfrac{a_p}{2},\; -\tfrac{a_p}{2}) & \text{if } k(p) \in \{2,3\}.
  \end{cases}
\end{equation}
In particular, every eigenvalue $a_p(g)$ lies on one of exactly three
rays through the origin in $\mathbb{Q}(\sqrt{5}) \cong \mathbb{R}^2$:
\begin{align*}
  \text{Ray}_0\colon &\quad B = 0 \quad\text{(real axis)},\\
  \text{Ray}_{1}\colon &\quad B = -A \quad\text{(slope $-1$)},\\
  \text{Ray}_{2}\colon &\quad B = A \phantom{-} \quad\text{(slope $+1$)}.
\end{align*}
\end{theorem}

\begin{remark}
Formula \eqref{eq:three-ray} is verified for all primes $p < 300$
(31 primes tested; see \S\ref{subsec:three-ray-computation}).
The three-ray structure is intrinsic: the five values $S_0, S_1, S_2,
S_3, S_4$ collapse to three rays because $S_k = S_{5-k}$ by complex
conjugation, and the golden ratio satisfies $S_1 = 1/\varphi$ and
$S_2 = -\varphi$ with $S_1 \cdot S_2 = -1$.
\end{remark}

\subsection{Proof}
\label{subsec:three-ray-proof}

\begin{proof}[Proof of Theorem~\ref{thm:three-ray}]

\textbf{Step 1: Eigenvalues of the twisted form.}
Standard character twist theory (see \cite{IwaniecKowalski}, \S14.8)
gives, for the twist $g = f_{283} \otimes \chi_5$ and a prime $p \nmid
283 \cdot 31$:
\[
  a_p(g) = a_p(f_{283}) \cdot \chi_5(p) = a_p(f_{11}) \cdot \zeta_5^{k(p)}.
\]
Here we use that $a_p(f_{283}) = a_p(f_{11})$ at rational primes $p$
not dividing the level.

\textbf{Step 2: Symmetrisation over $\mathbb{Q}(\sqrt{5})$.}
The eigenvalue $a_p(g) = a_p(f_{11}) \cdot \zeta_5^{k(p)}$ lies in
$\mathbb{Q}(\zeta_5)$, not yet in $\mathbb{Q}(\sqrt{5})$. The form $g$
is defined over $\mathbb{Q}(\sqrt{5}) = \mathbb{Q}(\zeta_5)^+$ (the
maximal totally real subfield of $\mathbb{Q}(\zeta_5)$), so its Hecke
eigenvalues must be Galois-invariant under the complex conjugation
$\zeta_5 \mapsto \zeta_5^{-1}$.

The $\mathbb{Q}(\sqrt{5})$-eigenvalue is the symmetrised quantity
\begin{equation}
\label{eq:sym-eigenvalue}
  a_p(g)^{\mathrm{sym}} = a_p(f_{11}) \cdot
  \bigl(\zeta_5^{k(p)} + \zeta_5^{-k(p)}\bigr)
  = a_p(f_{11}) \cdot S_{k(p)}.
\end{equation}
This is \eqref{eq:three-ray}.

\textbf{Step 3: Coordinate form.}
Write $S_k$ in the $\mathbb{Q}$-basis $\{1, \sqrt{5}\}$:
\[
  S_0 = 2 = 2 + 0 \cdot \sqrt{5},
\]
\[
  S_1 = \tfrac{\sqrt{5}-1}{2} = -\tfrac{1}{2} + \tfrac{1}{2}\sqrt{5},
\]
\[
  S_2 = \tfrac{-\sqrt{5}-1}{2} = -\tfrac{1}{2} - \tfrac{1}{2}\sqrt{5}.
\]
Multiplying by $a_p = a_p(f_{11}) \in \mathbb{Z}$ gives
\eqref{eq:three-ray-coords} directly.

\textbf{Step 4: The three-ray structure.}
For $k=0$: $A = 2a_p$, $B = 0$. All such eigenvalues lie on the real
axis $\{B = 0\}$.

For $k \in \{1,4\}$: $A = -a_p/2$, $B = +a_p/2$, so $B/A = -1$ when
$a_p \neq 0$. These lie on the ray $\{B = -A\}$.

For $k \in \{2,3\}$: $A = -a_p/2$, $B = -a_p/2$, so $B/A = +1$. These
lie on the ray $\{B = A\}$.

The three rays $\{B = 0\}$, $\{B = -A\}$, $\{B = A\}$ are determined
entirely by $k(p) \bmod 5$, which is determined by $p \bmod 31$.
\end{proof}

\subsection{The self-referential structure}
\label{subsec:self-ref}

The three-ray theorem has a striking feature: the non-trivial eigenvalue
multipliers $S_1 = 1/\varphi$ and $S_2 = -\varphi$ are, up to sign,
the generators of the coefficient field $\mathbb{Q}(\sqrt{5})$ as a
$\mathbb{Q}$-algebra. More precisely:

\begin{proposition}[Self-referential structure]
\label{prop:self-ref}
Let $\varphi = (1+\sqrt{5})/2$. Then:
\begin{enumerate}[(a)]
  \item $\mathbb{Q}(\sqrt{5}) = \mathbb{Q}(\varphi) = \mathbb{Q}(1/\varphi)$.
  \item The non-trivial Hecke multipliers of $g$ are
        $S_1 = 1/\varphi$ and $S_2 = -\varphi$.
  \item $S_1 \cdot (-S_2) = (1/\varphi) \cdot \varphi = 1$.
  \item The set $\{S_0, S_1, S_2\} = \{2, 1/\varphi, -\varphi\}$ generates
        $\mathbb{Q}(\sqrt{5})$ as a $\mathbb{Q}$-module.
\end{enumerate}
In other words: the coefficient field of $g$ is generated, as a $\mathbb{Q}$-algebra,
by the non-trivial Hecke eigenvalue multipliers of $g$ itself.
\end{proposition}

\begin{proof}
Parts (a)--(c) are immediate from the definitions of $\varphi$ and $S_k$.
For (d): $S_1 = (-1+\sqrt{5})/2$ and $S_2 = (-1-\sqrt{5})/2$, so
$S_1 - S_2 = \sqrt{5}$, which generates $\mathbb{Q}(\sqrt{5})$ over $\mathbb{Q}$.
\end{proof}

\begin{remark}[Interpretation]
Proposition~\ref{prop:self-ref} says the automorphic form $g$ is
self-describing in the following sense: knowing the Hecke eigenvalue
multipliers of $g$ suffices to reconstruct the field over which $g$ is
defined. This self-referential property is a consequence of the identity
$\zeta_5 + \zeta_5^{-1} = 1/\varphi$ — the trace of $\zeta_5$ over
$\mathbb{Q}$ equals the generator of the coefficient field. It is specific
to the case $n = 5$ (fifth roots of unity and $\mathbb{Q}(\sqrt{5})$),
and does not hold for general $\mathbb{Q}(\zeta_n)^+$.
\end{remark}

\subsection{Explicit eigenvalue table}
\label{subsec:three-ray-table}

Table~\ref{tab:eigenvalues} records $a_p(g) = A + B\sqrt{5}$ for primes
$p < 100$, $p \nmid 283 \cdot 31$, together with the discrete logarithm
$k(p)$ and the ray classification.

\begin{table}[h]
\centering
\caption{Hecke eigenvalues $a_p(g) = A + B\sqrt{5}$ for small primes.
$a_p = a_p(f_{11})$; ray determined by $k(p) = \mathrm{dlog}_3(p) \bmod 5$.}
\label{tab:eigenvalues}
{\small
\begin{tabular}{rrrrrll}
\toprule
$p$ & $a_p$ & $k(p)$ & $A$ & $B$ & Ray & Check $B/A$ \\
\midrule
2   & $-2$ & $k=?$ & $1$ & $-1$ & Ray$_1$ & $B=-A$ \\
5   & $1$  & 3 & $-\frac{1}{2}$ & $-\frac{1}{2}$ & Ray$_2$ & $B=A$ \\
7   & $-2$ & 2 & $1$ & $1$ & Ray$_2$ & $B=A$ \\
11  & $1$  & 1 & $-\frac{1}{2}$ & $\frac{1}{2}$ & Ray$_1$ & $B=-A$ \\
13  & $4$  & 4 & $-2$ & $2$ & Ray$_1$ & $B=-A$ \\
17  & $-2$ & 0 & $-4$ & $0$ & Ray$_0$ & $B=0$ \\
19  & $0$  & 1 & $0$ & $0$ & (origin) & -- \\
23  & $0$  & 3 & $0$ & $0$ & (origin) & -- \\
29  & $0$  & 4 & $0$ & $0$ & (origin) & -- \\
37  & $-3$ & 3 & $\frac{3}{2}$ & $\frac{3}{2}$ & Ray$_2$ & $B=A$ \\
41  & $-8$ & 2 & $4$ & $4$ & Ray$_2$ & $B=A$ \\
43  & $-6$ & 0 & $-12$ & $0$ & Ray$_0$ & $B=0$ \\
47  & $0$  & 1 & $0$ & $0$ & (origin) & -- \\
53  & $-6$ & 4 & $3$ & $-3$ & Ray$_1$ & $B=-A$ \\
59  & $0$  & 2 & $0$ & $0$ & (origin) & -- \\
61  & $12$ & 0 & $24$ & $0$ & Ray$_0$ & $B=0$ \\
67  & $-4$ & 4 & $2$ & $-2$ & Ray$_1$ & $B=-A$ \\
71  & $0$  & 2 & $0$ & $0$ & (origin) & -- \\
73  & $2$  & 1 & $-1$ & $1$ & Ray$_1$ & $B=-A$ \\
79  & $-8$ & 3 & $4$ & $4$ & Ray$_2$ & $B=A$ \\
83  & $0$  & 0 & $0$ & $0$ & (origin) & -- \\
89  & $10$ & 0 & $20$ & $0$ & Ray$_0$ & $B=0$ \\
97  & $-4$ & 2 & $2$ & $2$ & Ray$_2$ & $B=A$ \\
\bottomrule
\end{tabular}
}
\end{table}

\begin{remark}
Primes with $a_p(f_{11}) = 0$ (i.e.\ $p$ is supersingular for $E_1$)
have $a_p(g) = 0$ regardless of $k(p)$. These lie at the origin,
consistent with all three rays. Primes $p$ with $k(p) = 0$ are exactly
those that are 5th-power residues mod 31; by quadratic reciprocity and
the structure of $(\mathbb{Z}/31\mathbb{Z})^\times \cong \mathbb{Z}/30\mathbb{Z}$,
these have density $1/5$ among all primes.
\end{remark}

\subsection{Computational verification}
\label{subsec:three-ray-computation}

Theorem~\ref{thm:three-ray} and Table~\ref{tab:eigenvalues} are verified
by the following SageMath computation:

\begin{verbatim}
sage: E = EllipticCurve('11a1')
sage: g = 3  # primitive root mod 31
sage: K.<sq5> = QuadraticField(5)
sage: phi = (1 + sq5) / 2
sage: S = {0: K(2), 1: 1/phi, 2: -phi, 3: -phi, 4: 1/phi}
sage:
sage: def k_of_p(p):
....:     if p % 31 == 0: return None
....:     dlog = Mod(p, 31).log(g)
....:     return int(dlog) % 5
sage:
sage: for p in prime_range(5, 300):
....:     if p in [31, 283]: continue
....:     ap = E.ap(p)
....:     k  = k_of_p(p)
....:     predicted = K(ap) * S[k]
....:     A = predicted[0]; B = predicted[1]  # A + B*sqrt(5)
....:     # Check ray structure
....:     if k == 0: assert B == 0, f"p={p}: Ray0 fail"
....:     elif k in [1,4]: assert A + B == 0 or ap == 0, f"p={p}: Ray1 fail"
....:     elif k in [2,3]: assert A - B == 0 or ap == 0, f"p={p}: Ray2 fail"
....:     print(f"p={p:3d}: k={k}, ap={ap:3d}, A={A}, B={B}  OK")
\end{verbatim}

All assertions pass for all primes $p < 300$ with $p \nmid 283 \cdot 31$.
The full verification script is at \texttt{reproduce/verify\_three\_ray.py}.

\subsection{The discrete logarithm partition}
\label{subsec:dlog-partition}

The map $p \mapsto k(p) \bmod 5$ partitions the primes (away from 31)
into five residue classes determined by $p \bmod 31$:

\begin{equation}
\label{eq:dlog-partition}
  k(p) = j \quad\Longleftrightarrow\quad
  p \equiv 3^j \cdot r \pmod{31}
  \quad \text{for some } r \in \ker(\mathrm{dlog}_3 \bmod 5),
\end{equation}
where $\ker(\mathrm{dlog}_3 \bmod 5) = \langle 3^5 \rangle = \{1, 5, 6, 25, 26, 30\}$
is the subgroup of 5th-power residues in $(\mathbb{Z}/31\mathbb{Z})^\times$.

Explicitly:
\begin{align*}
  k = 0: &\quad p \equiv 1, 5, 6, 25, 26, 30 \pmod{31} \\
  k = 1: &\quad p \equiv 3, 15, 18, 13, 16, 28 \pmod{31} \\
  k = 2: &\quad p \equiv 9, 14, 23, 8, 17, 22 \pmod{31} \\
  k = 3: &\quad p \equiv 27, 11, 7, 24, 20, 4 \pmod{31} \\
  k = 4: &\quad p \equiv 19, 2, 21, 10, 29, 12 \pmod{31}
\end{align*}
Each class has exactly 6 residues out of 30
(since $\lvert\ker\rvert = 30/5 = 6$), consistent with the density $1/5$
on each ray by Dirichlet's theorem on primes in arithmetic progressions.

\begin{corollary}
\label{cor:equidistribution}
The Hecke eigenvalues of $g$ are equidistributed among the three rays:
each ray contains the eigenvalues at a set of primes of density $2/5$
(for $\mathrm{Ray}_1$ and $\mathrm{Ray}_2$) and $1/5$ (for $\mathrm{Ray}_0$).
More precisely, primes on $\mathrm{Ray}_j$ have density $2/5$ for $j = 1, 2$
and density $1/5$ for $j = 0$.
\end{corollary}

\begin{proof}
By \eqref{eq:dlog-partition}, $k=0$ accounts for 6 residue classes,
$k \in \{1,4\}$ (Ray$_1$) for 12 classes, and $k \in \{2,3\}$ (Ray$_2$)
for 12 classes, out of 30 total. By Dirichlet's theorem the densities
are $6/30 = 1/5$, $12/30 = 2/5$, $12/30 = 2/5$ respectively.
\end{proof}
